% MT045 / explanation — canonical editable content.
线性回归预测为\(w_1x_1+\cdots+w_dx_d+b\)。若\(w_j=0\)，第\(j\)个输入不会直接影响预测，因此“让一些权重恰好为0”能完成特征选择。



\begin{itemize}

\item \textbf{A正确。} 零系数直接删除相应特征的贡献。

\item \textbf{B错误。} Ridge惩罚平方和，通常让权重整体变小，不像LASSO那样主动产生稀疏解。

\item \textbf{C错误。} OMP（orthogonal matching pursuit）每轮挑一个与当前残差最相关的候选特征，再对已选集合重新拟合。它是在当前局面作贪心选择，没有穷举所有同样大小的特征集合，因此不保证全局最佳。

\item \textbf{D错误。} 即使特征集合相同，训练目标也可能不同。OLS只顾平方误差；Ridge加平方惩罚；LASSO加绝对值惩罚，学出的系数自然可能不同。

\item \textbf{E正确。} 权重平面上的一点\((w_1,w_2)\)就是一组候选系数。把权重绝对值的总和限制在上限\(t>0\)以内，即\(|w_1|+|w_2|\le t\)，可选范围便是菱形内部。边界尖角落在坐标轴上，这些位置有一个系数为0。下面用一个具体损失说明最优解怎样落到尖角。

\end{itemize}


补充教学目标为\(L(w_1,w_2)=\frac12[(w_1-1.6)^2+(w_2-0.4)^2]\)。没有约束时，最优点是\(u=(1.6,0.4)\)。把损失相同的参数点连起来，得到以它为中心的圆，这些圆叫损失等高线；圆越小，损失越低。


左图允许\(|w_1|+|w_2|\le1\)，右图允许\(w_1^2+w_2^2\le1\)。两者都把无约束最优点排除在外。把等高线从中心向外扩大，第一条接触允许范围的圆，对应能取到的最小损失。


\StudyImage{CS5489/exam-review/Midterm/assets/answer-l1-l2-geometry.png}


\emph{整理者补充的二维约束示意。横轴为权重\(w_1\)，纵轴为权重\(w_2\)；灰色区域是可选参数，叉号\(u\)是无约束最优点，虚线是首次接触可选范围的损失等高线，实心圆点为约束下的最优点。}


左图在菱形尖角\((1,0)\)接触，第二个系数恰好为0。右图离\(u\)最近的圆上点位于原点指向\(u\)的方向，把\(u\)除以自身长度便能缩到半径1：\((1.6,0.4)/\sqrt{1.6^2+0.4^2}\approx(0.970,0.243)\)，两个系数都非零。L1的轴上尖角有利于产生零系数；实际哪些系数为0仍取决于数据与惩罚强度。一维代数解释见\StudyRef{Answers}{mt019}{MT019}。


\textbf{出处：} 2023A期中 Q6，原卷第3页（5分）


\textbf{答案依据：} 2023A期中 Q6 随卷／配套参考答案；条件与解释经复核。
